\documentclass[aps,prl,reprint]{revtex4-2}
% Explicit, opt-in helpers for local typography previews.
% Choose the helper only after journal AND article type are known.
% Heading presence and wording remain author decisions.
% Pass heading words without terminal punctuation, except Note added. below.
% Never redefine \section globally with these helpers.
\newcommand{\prlrunin}[1]{%
  \par\noindent\hspace*{\parindent}%
  {\normalfont\itshape #1---}\ignorespaces}
\newcommand{\apsletterrunin}[1]{%
  \par\noindent\hspace*{\parindent}%
  {\normalfont\itshape #1.}\ \ignorespaces}
\newcommand{\apssecondaryrunin}[1]{%
  \par\noindent\hspace*{\parindent}%
  {\normalfont #1:}\quad\ignorespaces}
\newcommand{\apsdeepprlrunin}[1]{%
  \par\noindent\hspace*{\parindent}%
  {\normalfont\itshape #1.}\ \ignorespaces}

\begin{document}
\title{Synthetic layout examples: PRL with selected headings}
\author{Typography demonstration only}\noaffiliation
\maketitle
\noindent These synthetic layout examples contain no research findings. They illustrate chosen headings, not a required paper structure. This is a REVTeX preview, not a published paper or publisher-identical facsimile.

\prlrunin{Introduction}
The title and em dash are italic, without a full stop before the dash. Upright body text follows on the same line. The first line is indented. The name and presence of this heading are optional.

\prlrunin{Illustrative model}
A descriptive main-text heading follows the same typography. Actual manuscripts determine the number and names of their sections from their scientific argument.

\apssecondaryrunin{Global fit}
This optional secondary heading is roman. A colon and an em space separate it from the body text.

\apsdeepprlrunin{Boundary conditions}
A further subtopic uses an italic heading, a period, an ordinary word space, and then upright body text. It does not use the main-heading em dash.

\prlrunin{Conclusions}
This example chooses a plural ending. A paper could use Conclusion, Outlook, another descriptive title, or no closing heading.

% The period is intentional: Note added is a specific PRL exception.
\prlrunin{Note added.}
This typography-only sentence shows the specific period-plus-dash exception. It does not announce newly discovered research.

\begin{acknowledgments}
\prlrunin{Acknowledgments}
This is an illustrative placeholder. No contributor or funding source is being asserted. The native PRL environment has no printed title, so the explicit helper supplies one.
\end{acknowledgments}

\prlrunin{Data availability}
This is an illustrative placeholder, not a factual data-availability statement. Authors must describe their real data.
\end{document}
